\documentclass[10pt,a4paper]{amsart}
\usepackage[margin=19mm]{geometry}
\usepackage{amsmath,amssymb,mathtools}
\usepackage[scr=rsfs,cal=euler]{mathalfa}
\usepackage{xfrac}
\usepackage[unicode,hidelinks,bookmarks=false]{hyperref}
\newcommand{\ie}{i.e.\ }
\newcommand{\cf}{cf.\ }
\newcommand{\resp}{respectively\ }
\newcommand{\Subsection}{\subsection}
\providecommand{\nobreakdash}{\nobreak\hbox{-}}
\setlength{\emergencystretch}{2em}
\multlinegap0pt
\usepackage{bm}
\usepackage{bbm}
\usepackage{dsfont}
\usepackage{xspace}
\usepackage{cancel}
\usepackage{mathtools}
\usepackage{amsthm}
\makeatletter
\@ifundefined{prop}{\newtheorem{prop}{Proposition}}{}
\@ifundefined{theorem}{\newtheorem{theorem}{Theorem}}{}
\@ifundefined{corollary}{\newtheorem{corollary}[theorem]{Corollary}}{}
\@ifundefined{lemma}{\newtheorem{lemma}[theorem]{Lemma}}{}
\makeatother
\newcommand{\ddiv}{\mathrm{div}}
\newcommand{\dif}{\mathrm{d}}
\title[Asymptotic stability of planar viscous shock wave to three-d]{Asymptotic stability of planar viscous shock wave to three-dimensional relaxed compressible Navier-Stokes equations}
\author{Renyong Guan, Yuxi Hu}
\date{Source: arXiv:2511.07799v1 (CC BY 4.0)}
\begin{document}
\maketitle

\noindent\emph{Source and licence.} Asymptotic stability of planar viscous shock wave to three-dimensional relaxed compressible Navier-Stokes equations. Version arXiv:2511.07799v1 (CC BY 4.0), distributed under the Creative Commons Attribution 4.0 International licence, as recorded in the arXiv licence metadata for this exact version and shown on the abstract page https://arxiv.org/abs/2511.07799v1. Full rights notice: licenses/arxiv-2511.07799-CC-BY-4.0.txt.\par\smallskip

\noindent\textit{Editorial scope.} Counted unit: the Lemma whose source label is \texttt{lehs3} in the article (source0.tex lines 2923--2942, source label \texttt{lehs3}), delivered together with the article's own complete original proof of it (source0.tex lines 2943--3346, one \texttt{proof} environment of 404 lines). No mathematics is written, repaired or re-derived: statement and proof are the article's own text line for line. Required context retained uncounted: 1.11rh (lines 119--124); 1.12 (lines 126--128); 1.13 (lines 161--164); sccs (lines 176--179); 2.3 (lines 260--269); pvsw (lines 320--340); 3.1 (lines 392--401); th3.1 (lines 405--433); p1 (lines 461--484); sys1 (lines 689--711); le4.5 (lines 756--804); le0 (lines 1463--1487); legj (lines 1715--1744); sys3 (lines 1747--1779); coro111 (lines 2895--2920). Every definition, display and label of those lines is retained unchanged. Omitted from this package and not redistributed: 1--118, 125--125, 129--160, 165--175, 180--259, 270--319, 341--391, 402--404, 434--460, 485--688, 712--755, 805--1462, 1488--1714, 1745--1746, 1780--2894, 2921--2922, 3347--3592. Nothing inside a retained excerpt is omitted or altered: every retained body is delivered exactly as the article's lines are written, so no omission receipt of any kind accompanies this package. Citations: the retained excerpts carry no citation command at all, so no bibliography entry of the article is transcribed and no reference is redistributed. Reference adaptation: none was needed and none was made; every reference inside the delivered unit resolves inside it, so no unresolved reference or citation is produced. Printed slips kept literal: no printed word, symbol, hypothesis or number of the counted statement or of its proof was altered. Layout and class: the article's own class is \texttt{amsart} with packages including amssymb, amsmath, amsthm, mathrsfs, color, hyperref; the wrapper is the lane's pinned \texttt{amsart} editorial wrapper at 10pt on A4, which restates the theorem-like environments the retained text needs and adds the article's own macro definitions that the retained text needs (\texttt{\textbackslash ddiv}, \texttt{\textbackslash dif}), with extra wrapper packages (bm, bbm, dsfont, xspace, cancel, mathtools). The delivered numbering is the unit's own. Assets: no retained span contains a figure environment, a graphics-inclusion command, a PDF-page inclusion, a table or a TikZ picture; no image file is copied into this package, the package declares no asset dependency and no image file is redistributed with it. Licence: version arXiv:2511.07799v1 of this article is distributed under the Creative Commons Attribution 4.0 International licence, as recorded in the arXiv licence metadata for this exact version and shown on the abstract page https://arxiv.org/abs/2511.07799v1. A full rights notice is delivered at licenses/arxiv-2511.07799-CC-BY-4.0.txt. Characters: every retained excerpt is pure ASCII, so the delivered engine, XeLaTeX with its default Latin Modern text font, reports no missing character.
\begin{equation}\label{1.11rh}
\begin{cases}
-\sigma(\rho_+-\rho_-)+(\rho_+u_{1+}-\rho_-u_{1-})=0,\\
-\sigma(\rho_+u_{1+}-\rho_-u_{1-})+(\rho_+u_{1+}^2-\rho_-u_{1-}^2)+(p(\rho_+)-p(\rho_-))=0,
\end{cases}
\end{equation}

\begin{equation}\label{1.12}
\rho_->\rho_+, \qquad u_{1-}>u_{1+},
\end{equation}

\begin{align} \label{1.13}
(v, u, \Pi_1, \Pi_2)(0,x)=(v_0, u_0, \Pi_{10}, \Pi_{20})(x)
\rightarrow(v_{\pm}, u_{\pm}, \mathrm O_{3\times3}, 0) \quad (x_1\rightarrow\pm\infty),
\end{align}

\begin{equation}\label{sccs}
\tau\leq \min\{\inf\limits_{z\in[v_-,v_+]}\frac{\frac{4\mu}{3}+\lambda}{2|\sigma_\ast^2+p^{\prime}(z)|},
1\},
\end{equation}

\begin{equation}\label{2.3}
\begin{cases}
\rho^s\left(-\sigma v^s_{\xi_1}+ u_1^sv^s_{\xi_1}\right)=(u_1^s)_{\xi_1},\\
\rho^s\left(-\sigma( u_1^s)_{\xi_1}+u_1^s(u_1^s)_{\xi_1}\right)+p(v^s)_{\xi_1}=
                            (\Pi_{11}^s)_{\xi_1}+(\Pi_2^s)_{\xi_1},\\
\tau \rho^s \left(-\sigma(\Pi^s_{1})_{\xi_1}+ u_1^s(\Pi^s_{1})_{\xi_1}\right)+\Pi^s_{1}=
                          \text{diag}\{ \frac{4\mu}{3}(u_1^s)_{\xi_1}, -\frac{2\mu}{3}(u_1^s)_{\xi_1}, -\frac{2\mu}{3}(u_1^s)_{\xi_1}\},\\
\tau \rho^s \left(-\sigma(\Pi_2^s)_{\xi_1}+ u_1^s(\Pi^s_2)_{\xi_1}\right)+\Pi^s_2=\lambda(u_1^s)_{\xi_1},
\end{cases}
\end{equation}

\begin{lemma}\label{pvsw}
Under the assumption \eqref{sccs} for $\tau$, given any states $(v_-,u_{1-})$, $(v_+,u_{1+})$, and $\sigma_\ast>0$ satisfying R-H condition \eqref{1.11rh} and Lax  entropy condition \eqref{1.12}, there exists a positive constant $C$ independent of $\tau$ such that the following is true. The planar 2-viscous shock wave solutions $(v^s,u_{1}^s,\Pi_{11}^s,\Pi_2^s)(\xi_1)$ connecting $(v_-, u_{1-}, O_{3\times3}, 0)$ and $(v_+, u_{1+}, O_{3\times3}, 0)$ exists uniquely and satisfy
\begin{align*}
v^s_{\xi_1}>0,\quad
(u_1^s)_{\xi_1}=-\sigma_\ast v^s_{\xi_1}<0,
\end{align*}
and
\begin{align*}
&|v^s(\xi_1)-v_-|\leq C\delta e^{-C\delta|\xi_1|},\quad
|u_1^s(\xi_1)-u_{1-}|\leq C\delta e^{-C\delta|\xi_1|},\qquad  \forall \xi_1<0,\\
&|v^s(\xi_1)-v_+|\leq C\delta e^{-C\delta|\xi_1|},\quad
|u_1^s(\xi_1)-u_{1+}|\leq C\delta e^{-C\delta|\xi_1|},\qquad  \forall \xi_1>0,\\
&|(v^s_{\xi_1},u^s_{1\xi_1})|\le C\delta^2e^{-C\delta|\xi_1|},\quad
|\partial_{\xi_1}^{k}(v^s, u_1^s)| \le C\delta v^s_{\xi_1},\quad\quad\qquad\quad
\forall \xi_1\in\mathbb{R},\\
&|(\Pi^s_{11}, \Pi^s_2)|\le Cv^s_{\xi_1},\qquad\qquad
|\partial^{k-1}_{\xi_1}(\Pi^s_{11}, \Pi^s_2)|\le C\delta v^s_{\xi_1},
\qquad\quad\quad\forall\xi_1\in\mathbb{R},
\end{align*}
for $k=2,3,4,...$, where $\delta$ denote the strength of the shock wave as $\delta:=|p(v_-)-p(v_+)|\sim|v_--v_+|\sim|u_{1-}-u_{1+}|$.
\end{lemma}

\begin{equation}\label{3.1}
\begin{cases}
\rho\left(\partial_tv-\sigma\partial_{\xi_1}v+u\cdot\nabla_{\xi} v\right)=\ddiv_{\xi} u,\\
\rho\left(\partial_tu-\sigma\partial_{\xi_1}u+u\cdot\nabla_{\xi} u\right)+\nabla_{\xi}p(v)=\ddiv_{\xi}\Pi_1+\nabla_{\xi}\Pi_2,\\
\tau\rho\left(\partial_t\Pi_1-\sigma\partial_{\xi_1}\Pi_1+u\cdot\nabla_{\xi} \Pi_1\right)+\Pi_1
   =\mu\left(\nabla_{\xi} u+(\nabla_{\xi} u)^T-\frac{2}{3}\ddiv_{\xi} u \mathrm I_3\right),\\
\tau\rho\left(\partial_t\Pi_2-\sigma\partial_{\xi_1}\Pi_2+u\cdot\nabla_{\xi}\Pi_2\right)+\Pi_2
=\lambda\ddiv_{\xi} u.
\end{cases}
\end{equation}

\begin{theorem}\label{th3.1}
Let $\underline{v}$ and $\underline{u}_1$ be smooth monotone functions satisfying
\[
\underline{v}(x_1)=v_{\pm} \quad \underline{u}_1(x_1)=u_{1\pm} \quad for \quad \pm x_1\geq1.
\]
For any constants $M_0,M_1,\kappa_1,\kappa_2,\kappa_3,\kappa_4$ with $M_1>M_0>0$ and $\kappa_1>\kappa_2>\kappa_3>\kappa_4>0$, there exists a constant $T_0>0$ such that if
\[
\begin{aligned}
&\|v_0-\underline{v}\|_{H^3}+\|u_0-\underline{u}\|_{H^3}
+\sqrt{\tau}\|(\Pi_{10}, \Pi_{20})\|_{H^3}\leq M_0,\\
&0<\kappa_3\leq v_0(x)\leq\kappa_2,\qquad \forall x\in\Omega,
\end{aligned}
\]
where $\underline{u}=(\underline{u}_1, 0, 0)^T$, then the system \eqref{3.1} with initial condition \eqref{1.13} admits a unique solution $(v, u, \Pi_1, \Pi_2)$ on $[0,T_0]$ with
\[
\begin{aligned}
(v-\underline{v}, u-\underline{u}, \Pi_1, \Pi_2)\in C([0,T_0];H^3),
\end{aligned}
\]
and
\[
\|v-\underline{v}\|_{L^{\infty}(0,T_0;H^3)}+\|u-\underline{u}\|_{L^{\infty}(0,T_0;H^3)}
+\sqrt{\tau}\|(\Pi_1, \Pi_2)\|_{L^{\infty}(0,T_0;H^2)}\leq M_1.
\]
Moreover,
\[
\kappa_4\leq v(t,x)\leq\kappa_1,\qquad \forall(t,x)\in[0,T_0]\times\Omega.
\]
\end{theorem}

\begin{prop}\label{p1}
Let $(v, u, \Pi_1, \Pi_2)$ be local solutions given by Theorem \ref{th3.1} on $[0,T]$ for some $T>0$ and $( v^s,  u^s, \Pi^s_1, \Pi^s_{2})$ be defined in \eqref{2.3}. Then there exist positive constants $\delta_0,\varepsilon_1$ such that if the shock waves strength satisfy $\delta<\delta_0$ and
\begin{equation}\label{3.4}
\sup_{0\le t\le T} \|(v-(v^s)^{-X}, u-(u^s)^{-X}, \sqrt{\tau}(\Pi_1-(\Pi^s_1)^{-X}), \sqrt{\tau}(\Pi_2-(\Pi^s_2)^{-X}))\|_{H^3}\leq\varepsilon_1,
\end{equation}
then the following estimates hold
\begin{equation}\label{3.5}
\begin{aligned}
&\sup\limits_{t\in[0, T]}\|(v-(v^s)^{-X}, u-(u^s)^{-X}, \sqrt{\tau}(\Pi_1-(\Pi^s_1)^{-X}), \sqrt{\tau}(\Pi_2-(\Pi^s_2)^{-X}))\|_{H^3}\\
&\quad+\int_0^{T}\|\sqrt{|(v^s)^{-X}_{\xi_1}|}|v-(v^s)^{-X}|\|_{L^2}^2\dif t
+\int_0^{T}\|\left(\nabla_\xi(v-(v^s)^{-X}),\nabla_\xi (u-(u^s)^{-X})\right)\|^2_{H^2}\dif t
\\
&\quad+\delta\int_0^{T}|\dot X(t)|^2\dif t
+\int_0^{T}\|(\Pi_1-(\Pi_1^s)^{-X}, \Pi_2-(\Pi_2^s)^{-X})\|^2_{H^3}\dif t\\
&\leq C_0\left(\|(v_0-v^s, u_0-u^s, \sqrt{\tau}(\Pi_{10}-\Pi_1^s), \sqrt{\tau}(\Pi_{20}-\Pi_2^s))\|^2_{H^3}\right),
\end{aligned}
\end{equation}
where $C_0$ independent of $T$ and $\tau$.

In addition, 
\begin{equation}\label{3.6}
|\dot{X}(t)|\leq C_0\|(v-(v^s)^{-X})(t,\cdot)\|_{L^{\infty}},\qquad\forall t\leq T.
\end{equation}
\end{prop}

\begin{equation}\label{sys1}
\begin{cases}
\begin{aligned}
&\rho\partial_t\left(v-(v^s)^{-X}\right)-\sigma\rho\partial_{\xi_1}\left(v-(v^s)^{-X}\right)
        +\rho u\cdot\nabla_{\xi}\left(v-(v^s)^{-X}\right)-\rho\dot{X}(t)\partial_{\xi_1}(v^s)^{-X}
        +F\partial_{\xi_1}(v^s)^{-X}\\
 &\qquad=\ddiv_{\xi}\left(u-(u^s)^{-X}\right),\\
&\rho\partial_t\left(u-(u^s)^{-X}\right)-\sigma\rho\partial_{\xi_1}\left(u-(u^s)^{-X}\right)
        +\rho u\cdot\nabla_{\xi}\left(u-(u^s)^{-X}\right)-\rho\dot{X}(t)\partial_{\xi_1}(u^s)^{-X} +F\partial_{\xi_1}(u^s)^{-X}\\
&\qquad  +\nabla_{\xi}\left(p(v)-p((v^s)^{-X})\right)
               =\ddiv_{\xi}\left(\Pi_1-(\Pi^s)_1^{-X}\right)+\nabla_{\xi}\left(\Pi_2-(\Pi_2^s)^{-X}\right),\\
&\tau\rho\partial_t\left(\Pi_1-(\Pi_1^s)^{-X}\right)-\tau\sigma\rho\partial_{\xi_1}\left(\Pi_1-(\Pi_1^s)^{-X}\right)
+\tau\rho u\cdot\nabla_{\xi} \left(\Pi_1-(\Pi_1^s)^{-X}\right)-\tau\rho\dot{X}(t)\partial_{\xi_1}(\Pi_1^s)^{-X}\\
&\qquad +\tau F\partial_{\xi_1}(\Pi_1^s)^{-X}+\Pi_1-(\Pi_1^s)^{-X}\\
 &  =\mu\left(\nabla_{\xi}\left(u-(u^s)^{-X}\right)+\left(\nabla_{\xi}\left(u-(u^s)^{-X}\right)\right)^T
   -\frac{2}{3}\ddiv_{\xi}\left(u-(u^s)^{-X}\right) \mathrm{I}_3\right),\\
&\tau\rho\partial_t\left(\Pi_2-(\Pi_2^s)^{-X}\right)-\tau\sigma\rho\partial_{\xi_1}\left(\Pi_2-(\Pi_2^s)^{-X}\right)
+\tau\rho u\cdot\nabla_{\xi} \left(\Pi_2-(\Pi_2^s)^{-X}\right)-\tau\rho\dot{X}(t)\partial_{\xi_1}(\Pi_2^s)^{-X}\\
&+\tau F\partial_{\xi_1}(\Pi_2^s)^{-X}+\Pi_2-(\Pi_2^s)^{-X}
=\lambda\ddiv_{\xi} \left(u-(u^s)^{-X}\right),
\end{aligned}
\end{cases}
\end{equation}

\begin{lemma}\label{le4.5}
We have
\begin{equation}\label{4.5}
\frac{\dif}{\dif t}\int_{\mathbb{T}^2}\int_{\mathbb{R}}a^{-X}\rho\eta(t,\xi)\dif \xi_1\dif \xi^{\prime}=\dot{X}(t)Y(t)+J^{bad}(t)-J^{good}(t),
\end{equation}
where
\begin{align*}
Y(t):&=-\int_{\mathbb{T}^2}\int_{\mathbb{R}}a^{-X}_{\xi_1}\rho\eta(t,\xi)\dif \xi_1\dif \xi^{\prime}
-\int_{\mathbb{T}^2}\int_{\mathbb{R}}a^{-X}\rho p^{\prime}((v^s)^{-X})(v^s)^{-X}_{\xi_1}(v-(v^s)^{-X})\dif \xi_1\dif \xi^{\prime}\\
&+\int_{\mathbb{T}^2}\int_{\mathbb{R}}a^{-X}\rho \left(u-(u^s)^{-X}\right)(u^s)^{-X}_{\xi_1}\dif \xi_1\dif \xi^{\prime}\\
 &+\int_{\mathbb{T}^2}\int_{\mathbb{R}}a^{-X} \frac{\tau\rho}{2\mu}\left(\Pi_1-(\Pi_1^s)^{-X}\right):(\Pi_1^s)^{-X}_{\xi_1}\dif \xi_1\dif \xi^{\prime}\\
&+\int_{\mathbb{T}^2}\int_{\mathbb{R}}a^{-X} \frac{\tau\rho}{\lambda}\left(\Pi_2-(\Pi_2^s)^{-X}\right)(\Pi_2^s)^{-X}_{\xi_1}\dif \xi_1\dif \xi^{\prime},
\end{align*}
\begin{align*}
J^{bad}(t):=&\frac{1}{2\sigma_\ast}\int_{\mathbb{T}^2}\int_{\mathbb{R}}a^{-X}_{\xi_1}
|p(v)-p((v^s)^{-X})|^2\dif \xi_1\dif \xi^{\prime}
+\sigma_\ast\int_{\mathbb{T}^2}\int_{\mathbb{R}}a^{-X}p(v|(v^s)^{-X})(v^s)^{-X}_{\xi_1}\dif \xi_1\dif \xi^{\prime}\\
&+\frac{\delta}{\nu}\int_{\mathbb{T}^2}\int_{\mathbb{R}}\frac{a^{-X}}{\sigma_\ast v}a^{-X}_{\xi_1}\left(u_1-(u_1^s)^{-X}\right)^2\dif \xi_1\dif \xi^{\prime}\\
&-\int_{\mathbb{T}^2}\int_{\mathbb{R}}a^{-X}_{\xi_1}\left(u_1-(u_1^s)^{-X}\right)\left(\Pi_2-(\Pi_2^s)^{-X}\right)\dif \xi_1\dif \xi^{\prime}\\
&-\int_{\mathbb{T}^2}\int_{\mathbb{R}}a^{-X}_{\xi_1}\left(\left(u_1-(u^s_1)^{-X}\right)
\left(\Pi_{11}-(\Pi_{11}^s)^{-X}\right)+u_2\Pi_{21}+u_3\Pi_{31}
\right)\dif \xi_1\dif \xi^{\prime}\\
&+\int_{\mathbb{T}^2}\int_{\mathbb{R}}
a^{-X}\frac{1}{v}\left((\Pi_{11}^s)_{\xi_1}+(\Pi_{2}^s)_{\xi_1}\right)(v-(v^s)^{-X})\left(u_1-(u_1^s)^{-X}\right)
\dif \xi_1\dif \xi^{\prime}\\
&+\int_{\mathbb{T}^2}\int_{\mathbb{R}}a^{-X}\frac{1}{v\sigma_\ast}
\left((\Pi_{11}^s)_{\xi_1}+(\Pi_{2}^s)_{\xi_1}\right)\left(u_1-(u_1^s)^{-X}\right)^2
\dif \xi_1\dif \xi^{\prime}
\\
&+\int_{\mathbb{T}^2}\int_{\mathbb{R}}a^{-X}_{\xi_1}F\eta(t,\xi)\dif \xi_1\dif \xi^{\prime}
-\int_{\mathbb{T}^2}\int_{\mathbb{R}}a^{-X}\frac{\tau}{2\mu}F\left(\Pi_1-(\Pi_1^s)^{-X}\right) : (\Pi_1^s)^{-X}_{\xi_1}\dif \xi_1\dif \xi^{\prime}\\
&-\int_{\mathbb{T}^2}\int_{\mathbb{R}}a^{-X}\frac{\tau}{\lambda}F\left(\Pi_2-(\Pi_2^s)^{-X}\right)
(\Pi_2^s)^{-X}_{\xi_1}\dif \xi_1\dif \xi^{\prime}
:=\sum\limits_{i=1}^{10}B_i(t),
\end{align*}
and
\begin{align*}
J^{good}(t)&:=\sigma_\ast\int_{\mathbb{T}^2}\int_{\mathbb{R}}a^{-X}_{\xi_1} H(v|(v^s)^{-X})\dif \xi_1\dif \xi^{\prime}
+\sigma_\ast\int_{\mathbb{T}^2}\int_{\mathbb{R}}a^{-X}_{\xi_1} \frac{u_2^2+u_3^2}{2}\dif \xi_1\dif \xi^{\prime}\\
&+\frac{\sigma_\ast}{2}\int_{\mathbb{T}^2}\int_{\mathbb{R}}a^{-X}_{\xi_1}
\Big|u_1-(u_1^s)^{-X}-\frac{p(v)-p((v^s)^{-X})}{\sigma_\ast}\Big|^2\dif \xi_1\dif \xi^{\prime}\\
&-\int_{\mathbb{T}^2}\int_{\mathbb{R}}a^{-X}\frac{\sigma_\ast}{v}p^{\prime}((v^s)^{-X})(v^s)^{-X}_{\xi_1}(v-(v^s)^{-X})^2\dif \xi_1\dif \xi^{\prime}\\
&+\sigma_\ast\int_{\mathbb{T}^2}\int_{\mathbb{R}}a^{-X}_{\xi_1} \frac{\tau|\Pi_1-(\Pi_1^s)^{-X}|^2}{4\mu}\dif \xi_1\dif \xi^{\prime}
+\int_{\mathbb{T}^2}\int_{\mathbb{R}}a^{-X}\frac{|\Pi_1-(\Pi_1^s)^{-X}|^2}{2\mu}\dif \xi_1\dif \xi^{\prime}\\
&+\sigma_\ast\int_{\mathbb{T}^2}\int_{\mathbb{R}}a^{-X}_{\xi_1} \frac{\tau|\Pi_2-(\Pi_2^s)^{-X}|^2}{2\lambda}\dif \xi_1\dif \xi^{\prime}
+\int_{\mathbb{T}^2}\int_{\mathbb{R}}a^{-X}\frac{|\Pi_2-(\Pi_2^s)^{-X}|^2}{\lambda}\dif \xi_1\dif \xi^{\prime}\\
&:=\sum\limits_{i=1}^{8}G_i(t).
\end{align*}
\end{lemma}

\begin{lemma}\label{le0}
Under the hypotheses of Proposition \ref{p1}, there exists constant $C>0$ independent of $\tau, \nu, \delta, \varepsilon_1, T$, such that for $t\in[0, T]$, we have
\begin{equation}\label{ljgu1.1}
\begin{aligned}
%&\|v-(v^s)^{-X}\|_{L^2}^2+\|u-(u^s)^{-X}\|_{L^2}^2+\tau\|\Pi_1-(\Pi_1^s)^{-X}\|_{L^2}^2
%+\tau\|\Pi_2-(\Pi_2^s)^{-X}\|_{L^2}^2\\
& \int_{\mathbb{T}^2}\int_{\mathbb{R}}\rho\eta(t,\xi)\dif \xi_1\dif \xi^{\prime}
+\frac{\delta}{4M}\int_0^t|\dot X(t)|^2\dif t
+(1-C\varepsilon_1)\int_0^tG_2(t)\dif t
+(1-C(\frac{\delta}{\nu}+\delta+\varepsilon_1)\int_0^tG_3(t)\dif t\\
&\quad
+(C_1-C(\varepsilon_1+\sqrt\nu+\delta))\int_0^tG^s(t)\dif t
+(1-C\varepsilon_1)\int_0^t\|(\frac{\Pi_1-(\Pi_1^s)^{-X}}{2\mu}, \frac{\Pi_2-(\Pi_2^s)^{-X}}{\lambda})\|_{L^2}^2\dif t\\
&\le C\left(\|v_0-v^s\|_{L^2}^2+\|u_0-u^s\|_{L^2}^2+\tau\|\Pi_{10}-\Pi_1^s\|_{L^2}^2
+\tau\|\Pi_{20}-\Pi_2^s\|_{L^2}^2\right)\\
&\quad+(\frac{5\mu}{6}+\frac{5\lambda}{8})(1+\chi+C(\nu+\varepsilon_1))\int_0^t\int_{\mathbb{T}^2}\int_{\mathbb{R}}
|p^\prime(v)||\nabla(v-(v^s)^{-X})|^2
\dif \xi_1\dif \xi^{\prime}\dif t,
\end{aligned}
\end{equation}
where the terms $G_2(t),  G_{3}(t)$ are defined in Lemma \ref{le4.5}, the parameter $0<\chi\ll 1$ chosen to be sufficiently small later, and
\[
G^s(t)=\int_{\mathbb{T}^2}\int_{\mathbb{R}}(v^s)^{-X}_{\xi_1}|p(v)-p((v^s)^{-X})|^2\dif \xi_1\dif \xi^{\prime}.
\]
\end{lemma}

\begin{lemma}\label{legj}
Under the hypotheses of Proposition \ref{p1}, there exists constant $C>0$ independent of $\tau, \nu, \delta, \varepsilon_1, T$, such that for $t\in[0, T]$, we have
\begin{equation}\label{gjgj1.1}
\begin{aligned}
&\|\nabla_\xi(v-(v^s)^{-X})\|_{H^2}^2+\|\nabla_\xi(u-(u^s)^{-X})\|_{H^2}^2
+\tau\|\nabla_\xi(\Pi_1-(\Pi_1^s)^{-X})\|_{H^2}^2
+\tau\|\nabla_\xi(\Pi_2-(\Pi_2^s)^{-X})\|_{H^2}^2\\
&\quad
+\int_0^t\|(\nabla_\xi(\Pi_1-(\Pi_1^s)^{-X}), \nabla_\xi(\Pi_2-(\Pi_2^s)^{-X}))\|_{H^2}\dif t\\
&\le C\Big(\|\nabla_\xi(v_0-v^s)\|_{H^2}^2+\|\nabla_\xi(u_0-u^s)\|_{H^2}^2
+\tau\|\nabla_\xi(\Pi_{10}-\Pi_1^s)\|_{H^2}^2
+\tau\|\nabla_\xi(\Pi_{20}-\Pi_2^s)\|_{H^2}^2\Big)\\
%+C(\varepsilon_1+\delta)\int_0^t\|(\Pi_1-(\Pi_1^s)^{-X}, \Pi_2-(\Pi_2^s)^{-X})\|_{L^2}\dif z\\
&+C(\varepsilon_1+\delta)\int_0^t
\left(\|\left(\nabla_\xi(v-(v^s)^{-X}), \nabla_\xi(u-(u^s)^{-X})\right)\|_{H^2}^2
+\|(\Pi_1-(\Pi_1^s)^{-X}, \Pi_2-(\Pi_2^s)^{-X})\|_{L^2}
\right)
\dif t\\
&
+C\delta^2\int_0^t|\dot X(t)|^2\dif t
+C(\varepsilon_1+\delta)\int_0^t
G^s(t)
\dif t
+C\delta\int_0^t
G_3(t)
\dif t,
\end{aligned}
\end{equation}
where $G_3(t)$ and $G^s(t)$ are defined in Lemma \ref{le4.5} and Lemma \ref{le0}, respectively.
\end{lemma}

\begin{equation}\label{sys3}
\begin{cases}
&\partial_t\partial_\xi^\alpha\Phi-\sigma\partial_{\xi_1}\partial_\xi^\alpha\Phi
        + \partial_\xi^\alpha\left(u\cdot\nabla_\xi\Phi\right)
        -\dot{X}(t)\partial_{\xi_1}\partial_\xi^\alpha(v^s)^{-X}
 +\partial_\xi^\alpha\left(vF\partial_{\xi_1}(v^s)^{-X}\right)
 =\partial_\xi^\alpha\left(v\ddiv_\xi\Psi\right),\\
&\partial_t\partial_\xi^\alpha\Psi-\sigma\partial_{\xi_1}\partial_\xi^\alpha\Psi
        +\partial_\xi^\alpha\left( u\cdot\nabla_\xi\Psi\right)
        -\dot{X}(t)\partial_{\xi_1}\partial_\xi^\alpha(u^s)^{-X}
        +\partial_\xi^\alpha\left(vF\partial_{\xi_1}(u^s)^{-X}\right)\\
&\qquad +\partial_\xi^\alpha\left(v\nabla_\xi\left(p(v)-p((v^s)^{-X})\right)\right)
               =\partial_\xi^\alpha\left(v\ddiv_\xi Q_1\right)
               +\partial_\xi^\alpha\left(v\nabla_\xi Q_2\right),\\
&\tau\partial_t\partial_\xi^\alpha Q_1
-\tau\sigma\partial_{\xi_1}\partial_\xi^\alpha Q_1
+\tau\partial_\xi^\alpha\left( u\cdot\nabla_\xi Q_1\right)
-\tau\dot{X}(t)\partial_{\xi_1}\partial_\xi^\alpha(\Pi_1^s)^{-X}
+\tau\partial_\xi^\alpha\left(v F\partial_{\xi_1}(\Pi_1^s)^{-X}\right)
\\
 & \qquad+\partial_\xi^\alpha\left(vQ_1\right)
 =\mu\partial_\xi^\alpha\left(v\Big(\nabla_\xi\Psi+\left(\nabla_\xi\Psi\right)^T
   -\frac{2}{3}\ddiv_\xi\Psi \mathrm I_3\Big)\right),\\
&\tau\partial_t\partial_\xi^\alpha Q_2
-\tau\sigma\partial_{\xi_1}\partial_\xi^\alpha Q_2
+\tau\partial_\xi^\alpha\left( u\cdot\nabla_\xi Q_2\right)
-\tau\dot{X}(t)\partial_{\xi_1}\partial_\xi^\alpha(\Pi_2^s)^{-X}
+\tau \partial_\xi^\alpha\left(vF\partial_{\xi_1}(\Pi_2^s)^{-X}\right)
\\
&\qquad+\partial_\xi^\alpha\left(vQ_2\right)
=\lambda\partial_\xi^\alpha\left(v\ddiv_\xi \Psi\right).
\end{cases}
\end{equation}

\begin{corollary} \label{coro111}
	Under the hypotheses of Proposition \ref{p1}, there exists constant $C>0$ independent of $\tau, \nu, \delta, \varepsilon_1, T$, such that for $t\in[0, T]$, we have
	\begin{equation}\label{hsgjcoll}
		\begin{aligned}
			&\|v-(v^s)^{-X}\|_{L^2}^2+\|u-(u^s)^{-X}\|_{L^2}^2
			+\tau\|\Pi_1-(\Pi_1^s)^{-X}\|_{L^2}^2
			+\tau\|\Pi_2-(\Pi_2^s)^{-X}\|_{L^2}^2\\
			&\quad
			+\int_0^t\|(\Pi_1-(\Pi_1^s)^{-X}, \Pi_2-(\Pi_2^s)^{-X})\|_{L^2}\dif t
			+\int_0^t\|(\nabla_\xi(v-(v^s)^{-X}), \nabla_\xi(u-(u^s)^{-X}))\|_{L^2}\dif t\\
			%&\int_0^t\int_{\mathbb{T}^2}\int_{\mathbb{R}}|\nabla_\xi(v-(v^s)^{-X})|^2\dif \xi_1\dif \xi^{\prime}\dif t+\int_0^t\int_{\mathbb{T}^2}\int_{\mathbb{R}}|\nabla_\xi(u-(u^s)^{-X})|^2\dif \xi_1\dif \xi^{\prime}\dif t\\
			&\quad+\delta\int_0^t|\dot{X}(t)|^2
			\dif t+\int_0^t(G_2(t)+G_3(t)+G^s(t))
			\dif t\\
			&\le
			%\chi_1\|\sqrt\rho(u-(u^s)^{-X})\|_{L^2}^2
			%+\chi_2\|\sqrt\rho (\Pi_1-(\Pi_1^s)^{-X},
			%\|_{L^2}^2+\frac{\tau}{4}\|
			%\Pi_2-(\Pi_2^s)^{-X})\|_{L^2}^2\\
			C\|(v_{0}-v^s, u_1-u^s,
			%\|_{H^3}^2+\tau\|
			\sqrt{\tau}(\Pi_{10}-\Pi_1^s), \sqrt{\tau}(\Pi_{20}-\Pi_2^s))\|_{H^1}^2,
		\end{aligned}
	\end{equation}
	where $G_2(t), G_3(t)$ are defined in Lemma \ref{le4.5} and $G^s(t)$ is defined in Lemma \ref{le0}, respectively.
\end{corollary}

 \begin{lemma}\label{lehs3}
Under the hypotheses of Proposition \ref{p1}, there exists constant $C>0$ independent of $\tau, \nu, \delta, \varepsilon_1, T$, such that for $t\in[0, T]$, we have
\begin{equation}\label{hsgj1.3}
\begin{aligned}
&\int_0^t\left(\|\nabla^2_\xi(v-(v^s)^{-X})\|_{H^1}^2
 +\|\nabla^2_\xi(u-(u^s)^{-X})\|_{H^1}^2\right)\dif t\\&\le 	C\|(v_{0}-v^s, u_1-u^s,
%\|_{H^3}^2+\tau\|
\sqrt{\tau}(\Pi_{10}-\Pi_1^s), \sqrt{\tau}(\Pi_{20}-\Pi_2^s))\|_{H^3}^2.
%+C(\varepsilon_1+\delta)\int_0^t\|(\Pi_1-(\Pi_1^s)^{-X}, \Pi_2-(\Pi_2^s)^{-X})\|_{L^2}\dif z\\
%&\quad+C(\varepsilon_1+\delta)\int_0^t
%\left(\|\left(\nabla_\xi(v-(v^s)^{-X})\|_{L^2}^2+\| \nabla_\xi(u-(u^s)^{-X})\right)\|_{L^2}^2
%+G^s(t)\right)
%\dif t\\
%&\quad+C\delta^2\int_0^t|\dot X(t)|^2\dif t
%+C\delta\int_0^t
%G_3(t)
%\dif t,
\end{aligned}
\end{equation}
\end{lemma}

\begin{proof}
Multiplying $\eqref{sys3}_2$ by $\partial_\xi^\alpha\nabla_\xi\Phi$ for $|\alpha|=1, 2$, and integrating over $(0, t)\times\mathbb{R}\times\mathbb{T}^2$, we have
\begin{align*}
&-\int_0^t\int_{\mathbb{T}^2}\int_{\mathbb{R}}vp^\prime(v)|\partial_\xi^\alpha\nabla_\xi\Phi|^2
\dif \xi_1\dif \xi^{\prime}\dif t\\
&=\int_0^t\int_{\mathbb{T}^2}\int_{\mathbb{R}}\partial_t\partial_\xi^\alpha\Psi\cdot\partial_\xi^\alpha\nabla_\xi\Phi
\dif \xi_1\dif \xi^{\prime}\dif t
-\sigma\int_0^t\int_{\mathbb{T}^2}\int_{\mathbb{R}}\partial_{\xi_1}\partial_\xi^\alpha\Psi\cdot\partial_\xi^\alpha\nabla_\xi\Phi
\dif \xi_1\dif \xi^{\prime}\dif t\\
&\quad+\int_0^t\int_{\mathbb{T}^2}\int_{\mathbb{R}}\partial_\xi^\alpha\left( u\cdot\nabla_\xi\Psi\right)\cdot\partial_\xi^\alpha\nabla_\xi\Phi
\dif \xi_1\dif \xi^{\prime}\dif t
-\int_0^t\int_{\mathbb{T}^2}\int_{\mathbb{R}}\dot{X}(t)\partial_{\xi_1}\partial_\xi^\alpha(u^s)^{-X}\cdot\partial_\xi^\alpha\nabla_\xi\Phi
\dif \xi_1\dif \xi^{\prime}\dif t\\
&\quad+\int_0^t\int_{\mathbb{T}^2}\int_{\mathbb{R}}
\left(\partial_\xi^\alpha\left(v\nabla_\xi\left(p(v)-p((v^s)^{-X})\right)\right)-vp^\prime(v)\partial_\xi^\alpha\nabla_\xi\Phi\right)\cdot\partial_\xi^\alpha\nabla_\xi\Phi
\dif \xi_1\dif \xi^{\prime}\dif t\\
&\quad+\int_0^t\int_{\mathbb{T}^2}\int_{\mathbb{R}}\partial_\xi^\alpha\left(vF\partial_{\xi_1}(u^s)^{-X}\right)\cdot\partial_\xi^\alpha\nabla_\xi\Phi
\dif \xi_1\dif \xi^{\prime}\dif t
-\int_0^t\int_{\mathbb{T}^2}\int_{\mathbb{R}}\partial_\xi^\alpha\left(v\ddiv_\xi Q_1\right)\cdot\partial_\xi^\alpha\nabla_\xi\Phi
\dif \xi_1\dif \xi^{\prime}\dif t\\
&\quad
-\int_0^t\int_{\mathbb{T}^2}\int_{\mathbb{R}}\partial_\xi^\alpha\left(v\nabla_\xi Q_2\right)\cdot\partial_\xi^\alpha\nabla_\xi\Phi
\dif \xi_1\dif \xi^{\prime}\dif t=:\sum\limits_{i=1}\limits^8A_i.
\end{align*}
For $A_1$, applying equation $\eqref{sys3}_1$ and performing straightforward computation yields
\begin{align*}
A_1=&\int_{\mathbb{T}^2}\int_{\mathbb{R}}\left(\partial_\xi^\alpha\Psi\cdot\partial_\xi^\alpha\nabla_\xi\Phi
-\partial_\xi^\alpha\Psi_0\cdot\partial_\xi^\alpha\nabla_\xi\Phi_0\right)
\dif \xi_1\dif \xi^{\prime}
-\sigma\int_0^t\int_{\mathbb{T}^2}\int_{\mathbb{R}}
\partial_\xi^\alpha\Psi\cdot\partial_{\xi_1}\partial_\xi^\alpha\nabla_\xi\Phi
\dif \xi_1\dif \xi^{\prime}\dif t\\
&
+\int_0^t\int_{\mathbb{T}^2}\int_{\mathbb{R}}
\partial_\xi^\alpha\Psi\cdot\partial_\xi^\alpha\nabla_\xi\left(u\cdot\nabla_\xi\Phi\right)
\dif \xi_1\dif \xi^{\prime}\dif t
-\int_0^t\int_{\mathbb{T}^2}\int_{\mathbb{R}}\dot X(t)
\partial_\xi^\alpha\Psi\cdot\partial_{\xi_1}\partial_\xi^\alpha\nabla_\xi(v^s)^{-X}
\dif \xi_1\dif \xi^{\prime}\dif t\\
&
+\int_0^t\int_{\mathbb{T}^2}\int_{\mathbb{R}}
\partial_\xi^\alpha\Psi\cdot\partial_\xi^\alpha\nabla_\xi\left(vF\partial_{\xi_1}(v^s)^{-X}\right)
\dif \xi_1\dif \xi^{\prime}\dif t
-\int_0^t\int_{\mathbb{T}^2}\int_{\mathbb{R}}
\partial_\xi^\alpha\Psi\cdot\partial_\xi^\alpha\nabla_\xi\left(v\ddiv_\xi\Psi\right)
\dif \xi_1\dif \xi^{\prime}\dif t\\
=:&\sum\limits_{j=1}\limits^6A_{1j}.
\end{align*}
Using Young's inequality, we deduce that
\[
A_{11}\le \frac{1}{2}\left(\|\partial_\xi^\alpha\Psi\|_{L^2}^2+\|\partial_\xi^\alpha\nabla_\xi\Phi\|_{L^2}^2
+\|\partial_\xi^\alpha\Psi_0\|_{L^2}^2+\|\partial_\xi^\alpha\nabla_\xi\Phi_0\|_{L^2}^2\right).
\]
By integration by parts, we get
\[
A_{12}=\sigma\int_0^t\int_{\mathbb{T}^2}\int_{\mathbb{R}}
\partial_{\xi_1}\partial_\xi^\alpha\Psi\cdot\partial_\xi^\alpha\nabla_\xi\Phi
\dif \xi_1\dif \xi^{\prime}\dif t
=-A_2.
\]
Next, using similar estimates as in Lemma \ref{legj}, we obtain that
\begin{align*}
A_{13}\le C(\varepsilon_1+\delta)\int_0^t\left(\|\partial_\xi^\alpha\Psi\|_{H^1}^2+\|\nabla_\xi\Phi\|_{H^{|\alpha|}}^2\right)\dif t
+\int_0^t\int_{\mathbb{T}^2}\int_{\mathbb{R}}\partial_\xi^\alpha\Psi\cdot\left(u\cdot\partial_\xi^\alpha\nabla_\xi^2\Phi\right)
\dif \xi_1\dif \xi^{\prime}\dif t,
\end{align*}
\begin{align*}
A_{14}
\le C\delta^2\int_0^t|\dot X(t)|^2\dif t+C\delta\int_0^t\|\partial_\xi^\alpha\Psi\|_{L^2}^2\dif t,
\end{align*}
\begin{align*}
A_{15}\le C\delta\int_0^t\left(\|\nabla_\xi\Psi\|_{H^{|\alpha|}}^2+\|\nabla_\xi\Phi\|_{H^{|\alpha|}}^2+G_3(t)+G^s(t)\right)
\dif t,
\end{align*}
and
\begin{align*}
A_{16}
\le \int_0^t\int_{\mathbb{T}^2}\int_{\mathbb{R}}
v|\partial_\xi^\alpha\ddiv_\xi\Psi|^2
\dif \xi_1\dif \xi^{\prime}\dif t
+C(\delta+\varepsilon_1)\int_0^t\|\nabla_\xi\Psi\|_{H^{|\alpha|}}^2\dif t.
\end{align*}
Therefore, combining above estimates, we have
\begin{equation}\label{a1a2}
\begin{aligned}
A_1&+A_2\le \frac{1}{2}\left(\|\partial_\xi^\alpha\Psi\|_{L^2}^2+\|\partial_\xi^\alpha\nabla_\xi\Phi\|_{L^2}^2
+\|\partial_\xi^\alpha\Psi_0\|_{L^2}^2+\|\partial_\xi^\alpha\nabla_\xi\Phi_0\|_{L^2}^2\right)
+C\delta^2\int_0^t|\dot X(t)|^2\dif t\\
&
+C(\varepsilon_1+\delta)\int_0^t\left(\|\nabla_\xi\Psi\|_{H^{|\alpha|}}^2+\|\nabla_\xi\Phi\|_{H^{|\alpha|}}^2\right)\dif t
+\int_0^t\int_{\mathbb{T}^2}\int_{\mathbb{R}}
v|\partial_\xi^\alpha\ddiv_\xi\Psi|^2
\dif \xi_1\dif \xi^{\prime}\dif t
\\
&+C\delta\int_0^t\left(G_3(t)+G^s(t)\right)
\dif t
+\underbrace{\int_0^t\int_{\mathbb{T}^2}\int_{\mathbb{R}}\partial_\xi^\alpha\Psi\cdot\left(u\cdot\partial_\xi^\alpha\nabla_\xi^2\Phi\right)
\dif \xi_1\dif \xi^{\prime}\dif t}\limits_{=:A_{1\ast}}.
\end{aligned}
\end{equation}
Similarly, for $A_3$, we get
\begin{equation}\label{a3}
\begin{aligned}
A_3
\le C(\varepsilon_1+\delta)
\int_0^t\|(\nabla_\xi\Psi, \nabla_\xi\Phi)\|_{H^{|\alpha|}}^2\dif t
+\underbrace{\int_0^t\int_{\mathbb{T}^2}\int_{\mathbb{R}}\left( u\cdot\partial_\xi^\alpha\nabla_\xi\Psi\right)\cdot\partial_\xi^\alpha\nabla_\xi\Phi
\dif \xi_1\dif \xi^{\prime}\dif t}\limits_{=:A_{3\ast}},
\end{aligned}
\end{equation}
\begin{equation}\label{a4}
A_4\le C\delta^2\int_{0}^{t}|\dot{X}(t)|^2\dif t+C\delta\int_0^t\|\partial_\xi^\alpha\nabla_\xi\Phi\|_{L^2}^2\dif t,
\end{equation}
\begin{equation}\label{a5}
A_5\le C(\varepsilon_1+\delta)\int_0^t\left(G^s(z)+\|\nabla_\xi\Phi\|_{H^{|\alpha|}}^2\right)\dif t,
\end{equation}
\begin{equation}\label{a6}
\begin{aligned}
A_6\le C\delta\int_0^t\left(\|\nabla_\xi\Psi\|_{H^{|\alpha|-1}}^2+\|\nabla_\xi\Phi\|_{H^{|\alpha|}}^2+G_3(t)+G^s(t)\right)
\dif t,
\end{aligned}
\end{equation}
For $A_7$ and $A_8$, using Young's inequality, we get
\begin{equation}\label{a8}
\begin{aligned}
A_7+A_8\le-\frac{1}{8}\int_0^t\int_{\mathbb{T}^2}\int_{\mathbb{R}}vp^\prime(v)|\partial_\xi^\alpha\nabla_\xi\Phi|^2
\dif \xi_1\dif \xi^{\prime}\dif t+C\int_0^t\|(\nabla_\xi Q_1, \nabla_\xi Q_2)\|_{H^{|\alpha|}}^2\dif t
.
\end{aligned}
\end{equation}
%In a similar way, for $A_8$, we deduce that
%\begin{equation}\label{a8}%\begin{aligned}
%A_8\le-\frac{1}{16}\int_0^t\int_{\mathbb{T}^2}\int_{\mathbb{R}}vp^\prime(v)|\partial_\xi^\alpha\nabla_\xi\Phi|^2
%\dif \xi_1\dif \xi^{\prime}\dif t+C\int_0^t\|\nabla_\xi Q_2\|_{H^{|\alpha|}}^2\dif t
%.
%\end{aligned}
%\end{equation}
In addition, using Lemma \ref{pvsw}, we get
\begin{equation}\label{a9}
\begin{aligned}
A_{1\ast}+A_{3\ast}\le C(\delta+\varepsilon_1)\int_0^t\left(\|\partial_\xi^\alpha\Psi\|_{L^2}^2+\|\partial_\xi^\alpha\nabla_\xi\Phi\|_{L^2}^2\right)
\dif t.
\end{aligned}
\end{equation}

Therefore, combining \eqref{a1a2}-\eqref{a9} and choosing $\delta$ and $\varepsilon_1$ sufficiently small, we conclude that
\begin{equation}\label{gjhsgj3}
\begin{aligned}
&\frac{3}{4}\int_0^t\|\sqrt{|vp^\prime(v)}|\nabla_\xi^2\Phi\|_{H^1}^2
\dif \xi_1\dif \xi^{\prime}\dif t\\
&\le \frac{1}{2}\left(\|\nabla_\xi\Psi\|_{H^1}^2+\|\nabla_\xi^2\Phi\|_{H^1}^2
+\|\nabla_\xi\Psi_0\|_{H^1}^2+\|\nabla_\xi^2\Phi_0\|_{H^1}^2\right)
+C\delta^2\int_0^t|\dot X(t)|^2\dif t\\
&\quad
+C(\varepsilon_1+\delta)\int_0^t\left(\|\nabla_\xi\Psi\|_{H^2}^2+\|\nabla_\xi\Phi\|_{L^2}^2
+G^s(t)\right)\dif t
+C\delta\int_0^t
G_3(t)\dif t
\\
&\quad+\int_0^t\|
\sqrt{v}\nabla_\xi\ddiv_\xi\Psi\|_{H^1}^2
\dif t+C\int_0^t\left(\|\nabla_\xi Q_1\|_{H^2}^2+\|\nabla_\xi Q_2\|_{H^2}^2\right)
\dif t
\end{aligned}
\end{equation}

 Multiplying equation $\eqref{sys3}_3$ by $\partial_\xi^\alpha\nabla_\xi\Psi$ for $|\alpha|=1,2$, and integrating over $(0, t)\times\mathbb{R}\times\mathbb{T}^2$, we have
\begin{align*}
&\mu\int_0^t\int_{\mathbb{T}^2}\int_{\mathbb{R}}v|\partial_\xi^\alpha\nabla_\xi\Psi|^2\dif \xi_1\dif \xi^{\prime}\dif t
\\
=&\tau\int_0^t\int_{\mathbb{T}^2}\int_{\mathbb{R}}
\partial_t\partial_\xi^\alpha Q_1
:\partial_\xi^\alpha\nabla_\xi\Psi\dif \xi_1\dif \xi^{\prime}\dif t
-\tau\sigma\int_0^t\int_{\mathbb{T}^2}\int_{\mathbb{R}}
\partial_{\xi_1}\partial_\xi^\alpha Q_1
:\partial_\xi^\alpha\nabla_\xi\Psi\dif \xi_1\dif \xi^{\prime}\dif t\\
&+\tau\int_0^t\int_{\mathbb{T}^2}\int_{\mathbb{R}}
        \partial_\xi^\alpha ( u\cdot\nabla_\xi Q_1)
:\partial_\xi^\alpha\nabla_\xi\Psi\dif \xi_1\dif \xi^{\prime}\dif t
-\tau\int_0^t\int_{\mathbb{T}^2}\int_{\mathbb{R}}
     \dot{X}(t)\partial_{\xi_1}\partial_\xi^\alpha(\Pi_1^s)^{-X}
:\partial_\xi^\alpha\nabla_\xi\Psi\dif \xi_1\dif \xi^{\prime}\dif t\\
&+\tau\int_0^t\int_{\mathbb{T}^2}\int_{\mathbb{R}}
        \partial_\xi^\alpha(vF\partial_{\xi_1}(\Pi_1^s)^{-X})
:\partial_\xi^\alpha\nabla_\xi\Psi\dif \xi_1\dif \xi^{\prime}\dif t
+\int_0^t\int_{\mathbb{T}^2}\int_{\mathbb{R}}
    \partial_\xi^\alpha(v  Q_1)
:\partial_\xi^\alpha\nabla_\xi\Psi\dif \xi_1\dif \xi^{\prime}\dif t\\
&-\mu\int_0^t\int_{\mathbb{T}^2}\int_{\mathbb{R}}
\partial_\xi^\alpha\left(v\left(\nabla_\xi\Psi\right)^T
-\frac{2v}{3}\left(\ddiv_\xi\Psi I_3\right)\right):\partial_\xi^\alpha\nabla_\xi\Psi
\dif \xi_1\dif \xi^{\prime}\dif t=:\sum\limits_{i=1}\limits^7W_i.
\end{align*}
For $W_1$, by applying integration by parts and using $\eqref{sys1}_2$, we get
\begin{align*}
W_1
=&\tau\int_{\mathbb{T}^2}\int_{\mathbb{R}}
\left(\partial_\xi^\alpha Q_1
:\partial_\xi^\alpha\nabla_\xi\Psi
-\partial_\xi^\alpha Q_{10}
:\partial_\xi^\alpha\nabla_\xi\Psi_0\right)\dif \xi_1\dif \xi^{\prime}
-\tau\sigma\int_0^t\int_{\mathbb{T}^2}\int_{\mathbb{R}}
\partial_\xi^\alpha Q_1: \partial_{\xi_1}\partial_\xi^\alpha\nabla_\xi\Psi
\dif \xi_1\dif \xi^{\prime}\dif t\\
&+\tau\int_0^t\int_{\mathbb{T}^2}\int_{\mathbb{R}}
\partial_\xi^\alpha Q_1: \partial_\xi^\alpha\nabla_\xi(u\cdot\nabla_\xi\Psi)
        \dif \xi_1\dif \xi^{\prime}\dif t
+\tau\int_0^t\int_{\mathbb{T}^2}\int_{\mathbb{R}}
\partial_\xi^\alpha Q_1: (\dot{X}(t)\partial_{\xi_1}\partial_\xi^\alpha\nabla_\xi(u^s)^{-X})
    \dif \xi_1\dif \xi^{\prime}\dif t\\
&+\tau\int_0^t\int_{\mathbb{T}^2}\int_{\mathbb{R}}
      \partial_\xi^\alpha Q_1: \partial_\xi^\alpha\nabla_\xi(vF\partial_{\xi_1}(u^s)^{-X})
        \dif \xi_1\dif \xi^{\prime}\dif t\\
&-\tau\int_0^t\int_{\mathbb{T}^2}\int_{\mathbb{R}}
 \partial_\xi^\alpha Q_1: \partial_\xi^\alpha\nabla_\xi(v\nabla_\xi\left(p(v)-p((v^s)^{-X})\right))
    \dif \xi_1\dif \xi^{\prime}\dif t\\
    &+\tau\int_0^t\int_{\mathbb{T}^2}\int_{\mathbb{R}}
      \partial_\xi^\alpha Q_1: \partial_\xi^\alpha\nabla_\xi(v\ddiv_\xi Q_1)
        \dif \xi_1\dif \xi^{\prime}\dif t
-\tau\int_0^t\int_{\mathbb{T}^2}\int_{\mathbb{R}}
  \partial_\xi^\alpha Q_1: \partial_\xi^\alpha\nabla_\xi(v\nabla_\xi Q_2)
    \dif \xi_1\dif \xi^{\prime}\dif t\\
    &=:\sum\limits_{j=1}\limits^8W_{1j},
\end{align*}
Using the same method as for estimating $A_1$ and $A_5$, we imply that
\[
W_{11}\le \frac{1}{2}\left(\tau\|\partial_\xi^\alpha Q_{1}\|_{L^2}^2+\|\partial_\xi^\alpha\nabla_\xi\Psi\|_{L^2}^2
+\tau\|\partial_\xi^\alpha Q_{10}\|_{L^2}^2+\|\partial_\xi^\alpha\nabla_\xi\Psi_0\|_{L^2}^2\right),
\]
\[
W_{12}=\tau\sigma\int_0^t\int_{\mathbb{T}^2}\int_{\mathbb{R}}
\partial_{\xi_1}\partial_\xi^\alpha Q_1
:\partial_\xi^\alpha\nabla_\xi\Psi\dif \xi_1\dif \xi^{\prime}\dif t=-W_2,
\]
\begin{align*}
W_{13}&\le
 \frac{\mu}{24}\int_0^t\int_{\mathbb{T}^2}\int_{\mathbb{R}}v|\partial_\xi^\alpha\nabla_\xi\Psi|^2\dif \xi_1\dif \xi^{\prime}\dif t
+C\int_0^t\|\partial_\xi^\alpha Q_1\|_{H^1}^2\dif t
+C(\delta+\varepsilon_1)\int_0^t\|\nabla^k_\xi \Psi\|_{H^{|\alpha|-1}}^2\dif t,
\end{align*}
\[
W_{14}\le C\delta^2\int_0^t|\dot X(t)|^2\dif t+C\delta\int_0^t\|\partial_\xi^\alpha Q_1\|_{L^2}^2\dif t,
\]
\[
W_{15}\le C\delta\int_0^t\left(G_3(t)+G^s(t)+\|\partial_\xi^\alpha Q_1\|_{L^2}^2
+\|\nabla_\xi\Psi\|_{H^{|\alpha|}}^2+\|\nabla_\xi\Phi\|_{H^{|\alpha|}}^2\right)\dif t,
\]
\begin{align*}
W_{16}&\le \frac{-\mu}{200}\int_0^t\int_{\mathbb{T}^2}\int_{\mathbb{R}}vp^\prime(v)|\partial_\xi^\alpha\nabla_\xi\Phi|^2\dif \xi_1\dif \xi^{\prime}\dif t
    +C\int_0^t\|\partial_\xi^\alpha Q_1\|_{H^1}^2\dif t\\
    &\quad+C(\delta+\varepsilon_1)\int_0^t\left(G^s(z)+\|\nabla_\xi \Phi\|_{H^{|\alpha|-1}}^2\right)\dif t,
  \end{align*}
and
\[
 W_{17} \le C\int_0^t\|\nabla_\xi Q_1\|_{H^{|\alpha|}}^2\dif t,
 \quad
  W_{18} \le C\int_0^t\left(\|\partial_\xi^\alpha Q_1\|_{H^1}^2+\|\nabla_\xi Q_2\|_{H^{|\alpha|}}^2\right)\dif t.
\]
Thus, combining the above estimates, we derive that
\begin{equation}\label{w1}
\begin{aligned}
W_1+W_2\le& \frac{1}{2}\left(\tau\|\partial_\xi^\alpha Q_{1}\|_{L^2}^2+\|\partial_\xi^\alpha\nabla_\xi\Psi\|_{L^2}^2
+\tau\|\partial_\xi^\alpha Q_{10}\|_{L^2}^2+\|\partial_\xi^\alpha\nabla_\xi\Psi_0\|_{L^2}^2\right)
+C\delta^2\int_0^t|\dot X(t)|^2\dif t\\
&
+\frac{\mu}{24}\int_0^t\int_{\mathbb{T}^2}\int_{\mathbb{R}}v|\partial_\xi^\alpha\nabla_\xi\Psi|^2\dif \xi_1\dif \xi^{\prime}\dif t
+\frac{-\mu}{200}\int_0^t\int_{\mathbb{T}^2}\int_{\mathbb{R}}vp^\prime(v)|\partial_\xi^\alpha\nabla_\xi\Phi|^2\dif \xi_1\dif \xi^{\prime}\dif t\\
&+C(\delta+\varepsilon_1)\int_0^t\left(G_3(t)+G^s(t)+\|\nabla_\xi \Phi\|_{H^{|\alpha|}}^2+\|\nabla_\xi \Psi\|_{H^{|\alpha|}}^2\right)\dif t\\
&\quad+C\int_0^t\left(\|\nabla_\xi Q_1\|_{H^{|\alpha|}}^2+\|\nabla_\xi Q_2\|_{H^{\alpha}}^2\right)\dif t.
\end{aligned}
\end{equation}
Similarly, we obtain that
\begin{equation}\label{w3}
W_3\le \frac{\mu}{24}\int_0^t\int_{\mathbb{T}^2}\int_{\mathbb{R}}v|\partial_\xi^\alpha\nabla_\xi\Psi|^2\dif \xi_1\dif \xi^{\prime}\dif t
+C\int_0^t\|\partial_\xi^\alpha Q_1\|_{H^{1}}^2\dif t,
\end{equation}
\begin{equation}\label{w4}
W_4\le C\delta^2\int_0^t|\dot X(t)|^2\dif t+C\delta\int_0^t\|\partial_\xi^\alpha\nabla_\xi \Psi\|_{L^2}^2\dif t,
\end{equation}
\begin{equation}\label{w5}
W_5\le  C\delta\int_0^t\left(G_3(t)+G^s(t)
+\|\nabla_\xi\Psi\|_{H^{|\alpha|}}^2+\|\nabla_\xi\Phi\|_{H^{|\alpha|-1}}^2\right)\dif t,
\end{equation}
and
\begin{equation}\label{w6}
W_6\le \frac{\mu}{24}\int_0^t\int_{\mathbb{T}^2}\int_{\mathbb{R}}v|\partial_\xi^\alpha\nabla_\xi\Psi|^2\dif \xi_1\dif \xi^{\prime}\dif t
+C\int_0^t\|\nabla_\xi Q_1\|_{H^{|\alpha|-1}}^2\dif t.
\end{equation}
For $W_7$, first of all, straightforward computation yields
\begin{align*}
	&{\mu}\left(v\partial_\xi^\alpha\left(\nabla_\xi\Psi\right)^T
	-\frac{2v}{3}\partial_\xi^\alpha\left(\ddiv_\xi\Psi I_3\right)\right):\partial_\xi^\alpha\nabla_\xi\Psi
	=
	\mu\sum\limits_{i=1}\limits^3\sum\limits_{j=1}\limits^3
	\left(
	v \partial_\xi^\alpha\Psi_{i\xi_j}\cdot\partial_\xi^\alpha\Psi_{j\xi_i}
	-\frac{2v}{3} \partial_\xi^\alpha\Psi_{i\xi_i}\cdot\partial_\xi^\alpha\Psi_{j\xi_j}
	\right).
\end{align*}
By integration by parts, it holds
\begin{align*}
	&\int_0^t\int_{\mathbb{T}^2}\int_{\mathbb{R}}v \partial_\xi^\alpha\Psi_{i\xi_j}\cdot\partial_\xi^\alpha\Psi_{j\xi_i}
	\dif \xi_1\dif \xi^{\prime}\dif t\\
	&=
	-\int_0^t\int_{\mathbb{T}^2}\int_{\mathbb{R}}v_{\xi_i} \partial_\xi^\alpha\Psi_{i\xi_j}\cdot\partial_\xi^\alpha\Psi_{j}
	\dif \xi_1\dif \xi^{\prime}\dif t
	+\int_0^t\int_{\mathbb{T}^2}\int_{\mathbb{R}}v_{\xi_j} \partial_\xi^\alpha\Psi_{i\xi_i}\cdot\partial_\xi^\alpha\Psi_{j}
	\dif \xi_1\dif \xi^{\prime}\dif t\\
	&\quad+\int_0^t\int_{\mathbb{T}^2}\int_{\mathbb{R}}v \partial_\xi^\alpha\Psi_{i\xi_i}\cdot\partial_\xi^\alpha\Psi_{j\xi_j}
	\dif \xi_1\dif \xi^{\prime}\dif t.
\end{align*}
So,
\begin{align*}
	\sum\limits_{i=1}\limits^3\sum\limits_{j=1}\limits^3
	\int_0^t\int_{\mathbb{T}^2}\int_{\mathbb{R}}v \partial_\xi^\alpha\Psi_{i\xi_i}\cdot\partial_\xi^\alpha\Psi_{j\xi_j}
	\dif \xi_1\dif \xi^{\prime}\dif t
	=\int_0^t
	\|\sqrt{v} \partial_\xi^\alpha\ddiv_\xi\Psi\|_{L^2}^2
	\dif t
	,
\end{align*}
and
	\begin{align*}
		&\sum\limits_{i=1}\limits^3\sum\limits_{j=1}\limits^3
		\left(-\int_0^t\int_{\mathbb{T}^2}\int_{\mathbb{R}}v_{\xi_i} \partial_\xi^\alpha\Psi_{i\xi_j}\cdot\partial_\xi^\alpha\Psi_{j}
		\dif \xi_1\dif \xi^{\prime}\dif t
		+\int_0^t\int_{\mathbb{T}^2}\int_{\mathbb{R}}v_{\xi_j} \partial_\xi^\alpha\Psi_{i\xi_i}\cdot\partial_\xi^\alpha\Psi_{j}
		\dif \xi_1\dif \xi^{\prime}\dif t\right)\\
		&\le C(\delta+\varepsilon_1)
		\int_0^t \|\nabla_\xi\Psi\|^2_{H^{|\alpha|-1}}
		\dif t
		.
	\end{align*}
In addition,
	\begin{align*}
		&\mu\int_0^t\int_{\mathbb{T}^2}\int_{\mathbb{R}}
		\partial_\xi^\alpha\left(v\left(\nabla_\xi\Psi\right)^T
		-\frac{2v}{3}\left(\ddiv_\xi\Psi I_3\right)\right):\partial_\xi^\alpha\nabla_\xi\Psi
		\dif \xi_1\dif \xi^{\prime}\dif t\\
		&\quad-\mu\int_0^t\int_{\mathbb{T}^2}\int_{\mathbb{R}}
		\left(v\partial_\xi^\alpha\left(\nabla_\xi\Psi\right)^T
		-\frac{2v}{3}\partial_\xi^\alpha\left(\ddiv_\xi\Psi I_3\right)\right):\partial_\xi^\alpha\nabla_\xi\Psi
		\dif \xi_1\dif \xi^{\prime}\dif t\\
		&\le
		C(\delta+\varepsilon_1)
		\int_0^t \|\nabla_\xi\Psi\|^2_{H^{|\alpha|}}
		\dif t.
	\end{align*}
Thus, combining the above estimates, we derive that
\begin{equation}\label{w7}
W_7\le -\int_0^t
\|\sqrt{v} \partial_\xi^\alpha\ddiv_\xi\Psi\|_{L^2}^2
\dif t+C(\delta+\varepsilon_1)
\int_0^t \|\nabla_\xi\Psi\|^2_{H^{|\alpha|}}
\dif t.
\end{equation}

Therefore, combining \eqref{w1}-\eqref{w7} and using the smallness of $\delta$ and $\varepsilon_1$, we conclude that
\begin{equation}\label{gjhsgj4}
\begin{aligned}
&\frac{\mu}{2}\int_0^t\|\sqrt{v}\nabla^2_\xi\Psi\|_{H^1}^2\dif t
+\frac{\mu}{3}\int_0^t\|
\sqrt{v}\nabla_\xi\ddiv_\xi\Psi\|_{H^1}^2
\dif t\\
&\le
\frac{1}{2}\left(\tau\|\nabla_\xi Q_{1}\|_{L^2}^2+\|\nabla^2_\xi\Psi\|_{L^2}^2
+\tau\|\nabla_\xi Q_{10}\|_{L^2}^2+\|\nabla^2_\xi\Psi_0\|_{L^2}^2\right)
+C\delta^2\int_0^t|\dot X(t)|^2\dif t\\
&\quad
+\frac{-\mu}{50}\int_0^t\int_{\mathbb{T}^2}\int_{\mathbb{R}}\|\sqrt{|vp^\prime(v)|}\nabla^2_\xi\Phi\|_{H^1}^2\dif \xi_1\dif \xi^{\prime}\dif t
+C\int_0^t\left(\|\nabla_\xi Q_1\|_{H^{2}}^2+\|\nabla_\xi Q_2\|_{H^{2}}^2\right)\dif t\\
&\quad+C(\delta+\varepsilon_1)\int_0^t\left(G^s(t)+\|\nabla_\xi \Phi\|_{L^2}^2+\|\nabla_\xi \Psi\|_{L^2}^2\right)\dif t
+C\delta\int_0^tG_3(t)\dif t.
\end{aligned}
\end{equation}

Now, multiplying \eqref{gjhsgj3} and \eqref{gjhsgj4} by $\frac{\mu}{6}$ and $1$, respectively, and using the smallness of $\delta, \varepsilon_1$ together with Lemma \ref{legj}, we derive that
\begin{equation}\label{gjhsgj6}
\begin{aligned}
&\frac{\mu}{16}\int_0^t\|\sqrt{|vp^\prime(v)}|\nabla_\xi^2\Phi\|_{H^1}^2
\dif \xi_1\dif \xi^{\prime}\dif t
+\frac{\mu}{4}\int_0^t\|\sqrt{v}\nabla^2_\xi\Psi\|_{H^1}^2\dif t
+\frac{\mu}{6}\int_0^t\|
\sqrt{v}\nabla_\xi\ddiv_\xi\Psi\|_{H^1}^2
\dif t\\
&\le C\Big(\|\nabla_\xi\Phi_0\|_{H^2}^2+\|\nabla_\xi\Psi_0\|_{H^2}^2
+\tau\|\nabla_\xi Q_{10}\|_{H^2}^2
+\tau\|\nabla_\xi Q_{20}\|_{H^2}^2\Big)+C\delta^2\int_0^t|\dot X(t)|^2\dif t\\
%+C(\varepsilon_1+\delta)\int_0^t\|(\Pi_1-(\Pi_1^s)^{-X}, \Pi_2-(\Pi_2^s)^{-X})\|_{L^2}\dif z\\
&\quad+C(\varepsilon_1+\delta)\int_0^t
\left(\|\left(\nabla_\xi\Phi\|_{L^2}^2+\| \nabla_\xi\Psi\right)\|_{L^2}^2
\right)
\dif t
+C(\varepsilon_1+\delta)\int_0^t
G^s(t)
\dif t
+C\delta\int_0^t
G_3(t)
\dif t.
\end{aligned}
\end{equation}

Thus, combining Corollary \ref{coro111}, the proof of this is finished.
\end{proof}

\end{document}
